\chapter{Introducción}
\objetivo{Plantear el problema, la justificación, la pregunta y los objetivos, sin prometer un examen de sangre.}
\todojd{Redacta de 2 a 4 páginas. La guía de 2022 deja el planteamiento, la justificación y los objetivos dentro de la introducción, no como capítulos propios. Por eso este andamiaje no abre un capítulo para cada uno.}

\section{Campo y problema}
\begin{itemize}
\item El tejido publicado ya describe el tumor y la mucosa. \cite{vasaikar2019,zhang2014,tanaka2024}
\item El suero es otra matriz y otras personas. El depósito que lee esta tesis es PXD055218. \cite{rai2025}
\item CN no son personas sanas. HC se midió aparte y no decide la pregunta.
\end{itemize}
\todojd{Responde: qué se conoce del cáncer colorrectal en tejido y qué parte de ese problema aborda esta tesis. No localices las muestras en un segmento del colon.}

\section{Justificación}
\begin{itemize}
\item La equivalencia estadística no es identidad biológica. El cero de la clase diferente no dice que el tumor no cambie en el tejido.
\item No hay una cohorte nueva medida. Las redes de coexpresión no se calcularon.
\end{itemize}
\todojd{Escribe por qué la pregunta merece el trabajo y qué hueco no reclamas como prioridad.}

\section{Pregunta}
Qué parte del cambio del tumor se deja ver en la sangre cuando el contraste es otra enfermedad intestinal.
\todojd{Deja la pregunta en una frase y, en otra, sus límites: matrices separadas, sin examen clínico.}

\section{Objetivos}
\todojd{Redacta un objetivo general y los específicos en infinitivo. Una conclusión del capítulo 6 debe poder responder a cada específico.}
\begin{itemize}
\item General, a partir de la pregunta de arriba. No lo dejes como una promesa diagnóstica.
\item Específicos que el repo sí puede sostener: concordancia de signo en tejido (\TejidoNum\ de \TejidoDen); celda del suero (\SueroDif, \SueroEq, \SueroInd, de \SueroN); calidad de muestra y orden de adquisición; estado de H5 a H10.
\end{itemize}

\section{Alcance}
\begin{itemize}
\item Trabajo \textit{in silico} sobre datos públicos. No propone un examen clínico. \cite{pepe2001}
\item Certificado: \CertificadoEstado. Sello \texttt{71b310a3fcaebe9267fac971156df3aab491e009}, 2026-11-11T00:48:45Z. Etiqueta \texttt{v3.7-final} en \texttt{595118ae09ea6006c744c86b28cb0a10d42e0442}.
\end{itemize}
Citas sugeridas: \cite{vasaikar2019,tanaka2024,rai2025,garciafoncillas2026,pepe2001}.
Extensión orientativa: 2 a 4 páginas para todo el capítulo.
